\section{Background}

\DMGtodo{Should talk a bit about Hamenstaedt and Ledrappier work, copy some material of background of curves From curves to currents here, talk about definition of geodesic current in terms of cross-sections, etc.}

\begin{comment}
\DMGtodo[inline]{The following section contains material that could be relevant to include in intro or background section if we end up using anything for applications.}
\section{A bunch of random background and previous ideas}

\textbf{Cocycles}

A class of H\"older continuous function $\zeta\colon T^1 S \times \mathbb{R} \to \mathbb{R}$ satisfying $\zeta(v,s+t)=\zeta(\phi_s(v),t) + \zeta(v,s)$
for all $v \in T^1 S$ and all $s,t \in \mathbb{R}$.
Let $\mathcal{F}\colon T^1 S \to T^1 S$ be the flip $v \mapsto -v$. The flip of a cocycle $\zeta$ is the cocycle defined by $\mathcal{F}\zeta(v,t)=\zeta(\mathcal{F}\phi^t(v),t)$. Two cocycles $\zeta, \xi$ are cohomologous if $\zeta(v,t)-\xi(v,t)=\beta(\phi^t(v)) - \beta(v)$ for all $v \in T^1 S$ and all $t \in \mathbb{R}$.
$\zeta,\xi$ are cohomologous iff $\mathcal{F}\zeta,\mathcal{F}\xi$ are cohomologous. Thus, the flip map induces a map in cohomology classes.
We say $\zeta$ is quasi-invariant under the flip if $\zeta$ and $\mathcal{F}\zeta$ are cohomologous.

\textbf{Cross-ratio}

A H\"older continuous, $\pi_1(S)$-invariant function on $\partial \wt{S}^{[4]}$ (space of quadruples of pairwise distinct points of $\partial \wt{S}$) with properties

\begin{itemize}
\item $Cr$ is invariant under the diagonal action of $\pi_1(S)$.
\item $Cr(a,a',b,b')=Cr(a',a,b,b')^{-1}$
\item $Cr(a,a',b,b')=Cr(b,b',a,a')$
\item $Cr(a,a',b,b')Cr(a',a'',b,b')=Cr(a,a'',b,b')$
\end{itemize}



\textbf{Periods}

\textbf{Periods for cross-ratio}
Given a homotopy class of a closed curve $c$ in $S$, we can consider a representative $\gamma$.
Let $\gamma$ be the Deck transformation associated to $\gamma$, and let $\gamma^-$ and $\gamma^+$ be its repelling and attracting points respectively.
The period $l(\gamma)=\log Cr(\gamma^+,\gamma^-,\gamma \cdot a, a)$ for any $a \notin \{\gamma^-,\gamma^+\}$.
This definition does not depend on the conjugacy class of the element $\gamma \in \pi_1(S)$, nor on the choice of point $a$. \\

\textbf{Periods for geodesic currents}

Let $c$ be a closed curve in $S$, lift it to the universal cover to some $\gamma$, choosing a hyperbolic metric. Let $G_{\gamma}$ be the set of geodesics that intersect $\gamma$ and have one endpoint in the segment $(a,\gamma a)$. $\nu(G_{\gamma})$ is the same as the measure of the set of geodesics that intersect the segment $(q,\gamma q)$ tranversally, for $q$ in the axis.




\textbf{Signed geodesic current}

A geodesic current for $S$ is a $\pi_1(S)$-invariant locally finite Borel measure on $\partial \wt{S} \times \partial \wt{S} - \Delta$ which in addition is invariant under the natural involution exchanging the two factors. Geodesic currents naturally correspond to finite $\phi^t$ invariant transverse Borel measures on $T^1M$ invariant under the flip map.
A signed geodesic current is a difference of two geodesic currents. \\

Otal proves that given a family of periods there is at most one cross ratio with those periods (essentially, density of endpoints of hyperbolic elements).
What are the properties a period must satisfy?

\textbf{From flip invariant Holder class to cross-ratio}

Let $\alpha$ be a cocycle representing the class. Let $x \in \wt{S}$ be any point and let $(,)_{\alpha}$ be the symmetric product on distinct points in $T^1_x \wt{S}$. For $a \neq b \in \partial \wt{S}$, write $(a,b)=(v,w)_{\alpha}$ where $v,w \in T^1 \wt{S}$ are such that $\pi(v)=a$ and $\pi(w)=b$. Then $(,)$ is a symmetric product on pairs of distinct points in $\partial \wt{S}$ depending on the choice of $x$ and $\alpha$. For a quadruple $(a,a',b,b')$ of distinct points in $\partial \wt{S}$ define $[a,a',b,b']_{\alpha}=(a,b)+(a',b')-(a,b')-(a',b')$. Then one checks that $[]_{\alpha}$ does not depend on the choice of $x$ or the representative $\alpha$. And $Cr=e^{[]_{\alpha}}$ is the cross-ratio associated to the class of $\alpha$.

\textbf{From cross-ratio to Holder cocycle}

One defines a cocycle for the frame flow in the frame bundle, then integrates over fibers to obtain a Holder cocycle.\\

Finally one checks that one is inverse of the other (if one closes the cycle obtains a cohomologous cocycle). \\

\textbf{From geodesic current to cross-ratio}

Let $\mu$ be a geodesic current. Let $(a,b,c,d)$ an ordered tuple of pairwise distinct points. Define $[a,b,c,d]=\mu([a,b] \times [c,d])$. \\
(To check Holder property need to use $\mu$ is Gibbs current and thus comes from a Holder function...?) \\
To check positivity one uses:

\begin{lemma}
Let $Cr$ be a cross-ratio on $\partial \wt{S}$. Then $Cr(a,b,c,d)+Cr(b,c,d,a) \to 1$ as $a \to b$, locally uniformly in $(c,d)$ if and only if $Cr$ is positive.
\end{lemma}

\textbf{From cross-ratio to intersection geodesic current}

For a generalized cross-ratio there is a unique signed current $\mu$. Define $\mu([a,b]\times[c,d])=\log Cr(a,b,c,d)$. Additivity of cross ratio gives additivity of measure. If the cross ratio is positive then it gives a geodesic current. \\

\textbf{From cocycle to Gibbs current}

Let $\mathcal{M}$ be the space of all $\phi^t$-invariant probability measures on $T^1 S$.

The entropy of a Borel set $A$ w.r.to to probability measure $\nu$ on $T^1S$ is $h_{\nu}(A):=-\nu(A)\log \nu(A)$. \\
The entropy of $\nu$ is the entropy of the whole space $T^1 S$. Let $\mathcal{S}$ be the $\sigma$-algebra of all measurable sets. The entropy of $\nu$ is

$h_{\nu} = \sup \{ \sum_{A\mathcal{F}} h_{\nu}(A) \mathcal{F} \subset \mathcal{S} | \mathcal{F} \mbox{ finite } , X=\sqcup_{A \in \mathcal{S}} A\}$. \\

For every H\"older continuous function $f\colon T^1 S \to \mathbb{R}$ the pressure $p(f):=\sup \{ h_{\nu} - \int f d\nu | \nu \in \mathcal{M} \}$.

The Gibbs equilibrium state of $f$ is the (unique) probability measure $\eta_f \in \mathcal{M}$ which realizes this supremum.

Any H\"older function $f\colon T^1 S \to \mathbb{R}$ gives rise to a H\"older cocycle by 

$$\xi_f(v,t):=\int_0^{t} f (\phi^s(v)ds.$$

We call $f$ quasi-invariant under the flip if for any periodic $v \in T^1 S$ of period $\tau$, $\xi_f(v,\tau)=\xi_f(-v,\tau)$. \\

For a H\"older cocycle, we can find some H\"older function $f\colon T^1 S \to \mathbb{R}$. Two H\"older $f,g$ functions give the same cocycle iff $\int fd\mu = \int g d\mu$ for all $\mu \in \mathcal{M}$.\\

Otal proves that cross-ratio is determined given by periods. \\

\subsection{Reconstructing currents from periods}

Otal's argument uses density of geodesic in tangent bundle to construct current in some rectangles. Idea: extend to all sigma-algebra using Caratheodory's theorem. Problems:

\begin{itemize}
\item Is it enough to reconstruct the current given only on this family of rectangles?
\item Continuity at endpoints.
\end{itemize}

\subsection{Reconstructing current}

Another approach is to try to define a positive functional on the space of compact continuous functions on the space of geodesic currents. \\
It also has to have the $\pi_1(S)$-invariant property.\\
The difficult point will be continuity of the functional. \\


Feng Luo characterizes functions that are intersection numbers on simple curves (ref).
\end{comment}

\subsection{Action of $\pi_1(S,*)$ on lifts}
\label{subsec:actionsfundamental}

There are two actions of $\pi_1(S,*)$ on lifts of $*$ to $\tilde{S}=\{ \mbox{paths from } * \}/\sim$:

\begin{enumerate}
    \item Left action: precomposition with $\gamma$ via Deck transformations, isometries of $\mathbb{H}^2$. The isometry extends to a homeomorphism of the boundary at infinity $S^1_{\infty}$. This action is minimal, in the sense that any orbit $\pi_1(S,*).p$, for $p \in S^1_{\infty}$, is a dense subset of $S^1_{\infty}$.
    \item Right action: post-composition of paths. For example, $a \in \pi_1(S,*)$ acts on $x \in \pi_1(S,*)$, as $ax$, where this is the path first running through $x$ and then through $a$. 
\end{enumerate}

%\begin{definition}
%  A homogeneous function on geodesic currents is \emph{convex} if its
%  restriction to curves is convex by Definition~\ref{def:convex}.
%\end{definition}

%\begin{definition}
%  \emph{Pants coordinates} on $\mathcal{ML}(S)$ are Dehn-Thurston
%  coordinates with respect to a marked pair-of-pants decomposition,
%  consisting of length and twist coordinates for each pants curve. The
%  twist parameters require more specification, since there are several
%  variations in the literature; see Section XXX for our version.
%\end{definition}

%\begin{theorem}
%  Let $f \colon \mathcal{GC}(S) \to \R_{\ge 0}$ be a convex
%  functional. Then the restriction of $f$ to measured laminations is
%  convex with respect to pants coordinates for any pair-of-pants
%  decomposition.
%\end{theorem}

%\subsection{Outline}
%Extending functions to ML: \\
%- Convexity argument, natural coordinates? \\
%- Examples:
%\begin{enumerate}
%\item Hyperbolic length (Thurston)
%\item Intersection number (Thurston)
%\item Extremal length (Kerckhoff)
%\end{enumerate}

%Extending functions to GC:
%\begin{enumerate}
%\item Hyperbolic length (Bonahon)
%\item Intersection number (Bonahon)
%\item Stable length (Bonahon, Erlandsson-Parlier-Souto)
%\item Generalized translation lengths (Martone-Zhang)
%\end{enumerate}

%In this paper we provide an argument that generalizes the previous extensions to a broad family of functions.

%Applications:

%\begin{enumerate}
%\item Counting: Rafi-Souto, Erlandsson-Parlier-Souto)
%\item Extending extremal length to geodesic currents, relevant in study of %rational maps.
%\item Others? (conversations with Wright).
%\end{enumerate}

%Inequalities:
%\[C_1 l(c) \leq f(c) \leq C_2 l(c).\]

%For general $f$, the right inequality follows from argument decomposing curve into generators (in the additive case).
%Question: can we also obtain it in the non-additive case, for general $f$? \\ Explain the case $f$ is extremal length. Consider the case of the torus. \\
%Right inequality follows from collar lemma, and other by considering embedded annulus, and looking at the core curve. \\
%These inequalities are interesting because: (1) they are satisfied by all Liouville currents (see paper by Burger et al on structure theorem for geodesic currents), and (2) the left inequality is usually a part of the definition of Anosov representation, where $f$ plays the role of the generalized translation length (see Martone-Zhang).

\subsection{Background on measure theory}

Let $X$ be a set.

\begin{definition}[Semi-ring, ring, $\sigma$-ring]
A family $\mathcal{S}$ of subsets of a set $X$ is called a \emph{semiring} if it contains the empty set, $A \cap B \in \mathcal{S}$ for all $A,B \in \mathcal{S}$ and, for every pair of sets $A,B \in \mathcal{S}$ with $B A \subset B$, the set $B \backslash A$ is a disjoint union of finitely many sets in $\mathcal{S}$.
A family $\mathcal{R}$ of subsets of a set $X$ is called a \emph{ring} if it contains the empty set and the sets $A \cap B$, $A \cup B$ and $A \backslash B$ belong to $\mathcal{R}$ for all $A,B \in \mathcal{R}$. 
The \emph{$\sigma$-ring} of $\mathcal{R}$ consists of the set of countable unions of elements of $\mathcal{R}$.
\end{definition}

\begin{definition}[semialgebra, algebra, $\sigma$-algebra]
If in the above definitions one assumes $X$ is also an element of the family, one gets, respectively, the definitions of \emph{semialgebra},\emph{algebra} and $\sigma$-algebra.
\end{definition}

\begin{lemma}
For any semiring $\mathcal{S}$, the collections of all finite unions of sets in $\mathcal{S}$ forms a ring $\mathcal{R}_{\mathcal{S}}$, the \emph{ring generated by $\mathcal{S}$}. Every set in $\mathcal{R}$ is a finite union of pairwise disjoint sets in $\mathcal{S}$.
\end{lemma}
\begin{proof}
See \cite[1.2.14~Lemma]{Bo07:MeasureTheory}.
\end{proof}

\begin{definition}
A real-valued set function $\mu$ defined on a class of sets $\mathcal{A}$ is called \emph{finitely additive} if 

\[
\mu\left(\bigcup_{i=1}^n A_i \right) = \sum_{i=1}^n \mu(A_i)
\]
for all $n$ and all disjoint sets $A_1,\cdots,A_n \in \mathcal{A}$ such that $\bigcup_{i=1}^n A_i \in \mathcal{A}$.
$\mu$ is called \emph{countably subadditive} if 

\[
\mu\left(\bigcup_{i=1}^{\infty} A_i \right) = \sum_{i=1}^{\infty} \mu(A_i)
\]
for all pairwise disjoint sets $A_i \mathcal{A}$ such that $\bigcup_{i=1}^{\infty} A_i \in \mathcal{A}$.
We say $\mu$ is called \emph{countably additive} if the above inequality is replaced by an equality.
\end{definition}

\begin{proposition}
A finitely additive function $\mu$ on a semiring $\mathcal{S}$ extends uniquely to a finitely additive function on $\mathcal{R}_{S}$
\label{prop:addsemitoring}
\end{proposition}
\begin{proof}
See \cite[1.3.10~Proposition]{Bo07:MeasureTheory}.
\end{proof}


\begin{definition}
  Fix a ring $\mathcal{R}$. Let $\mathcal{F}$ be a family of nested sequences
  of elements of $\mathcal{R}$. Then $\mathcal{F}$ is \emph{final} if,
  for any sequence $(R_i)_{i=0}^{\infty}$ of nested elements of
  $\mathcal{R}$ with $\bigcap_i R_i = \emptyset$, there is a sequence
  $(S_j) \in \mathcal{F}$ so that $R_i$ is eventually contained in
  $S_j$, in the sense that for each $j$ there is an $N$ so that for
  all $i \ge N$ we have $R_i \subset S_j$.
\end{definition}

To show countable additivity, suffices to show that there is a final
family on which measures go to zero.


\begin{proposition}
Let $\mu$ be a finitely additive function defined on a ring $\mathcal{R}$. Let $\mathcal{F}$ be a final family of $\mathcal{R}$ so that for any $(S_j)_{j=1}^{\infty} \in \mathcal{F}$, $\lim_j \mu(S_j) \to 0$. Then $\mu$ is countably subadditive.
\label{prop:finalfamcount}
\end{proposition}
\begin{proof}
Let $A_n$ be a sequence of pair-wise disjoint sets in $\mathcal{R}$ whose union $A$ is also in $\mathcal{R}$.
Then consider $B_n =  A - \cup_{i=1}^n A_i$ . The sequence $B_n$ is a nested decreasing sequence of elements in $\mathcal{R}$ with empty intersection. Since $\mathcal{F}$ is a final family, for every $j$ we can find $N$ so that for every $i \geq N$, we have $B_i \subset S_j$.
On the one hand, note that $\mu(B_i) \to 0$, since $\mu(S_j) \to 0$ by definition of final family and $\mu$ being monotone.
On the other hand, \[
\mu(A - \cup_{i=1}^n A_i )= \mu(A) - \mu( \cup_{i=1}^n A_i )
\]
by finite additivity of $\mu$, which proves what we wanted.
\end{proof}


\begin{definition}\DMGtodo{Is this same as outer regular in Radon definition? What about inner regular?}
A set function $\mu$ defined on the class of all subsets of a set $X$ and taking values in $[0,\infty]$ is called an \emph{outer measure} if

\begin{enumerate}
    \item $\mu(\emptyset)=0$
    \item $\mu(A) \leq \mu(B)$ whenever $A \subset B$, i.e., is monotone
    \item $\mu$ is countably subadditive. 
\end{enumerate}
A set $A \subset X$ is called $\mu$-measurable if, for every $E \subset X$, one has the equality

\[
\mu(E \cap A) + \mu(E \backslash A) = \mu(E)
\]
The class of all $\mu$-measurable sets is denoted by $\mathcal{M}_{\mu}$.
An outer measure is called \emph{regular} if, for every set $A \subset X$, there exist a $\mu$-measurable set $B$ so that $A \subset B$ and $\mu(A)=\mu(B)$.
\end{definition}

\begin{theorem}[Caratheodory extension theorem]
Given a ring $\mathcal{R}$ and a finitely additive and countably additive function $\mu$ on it.
Let 
\[
\mu^*(A) \coloneqq \inf \{ \sum_{i=1}^n \mu(A_n) | A_n \in \mathcal{R}, A \subset \cup_{i=1}^{\infty} X_i\},
\]
where in the case of absence of such sets $A_n$ we set $\mu^*(A) \coloneqq \infty$. Then

\begin{enumerate}
    \item  $\mu^*$ is an outer regular measure
    \item All the sets in $\mathcal{R}$ are $\mu$-measurable
    \item $\mu^*$ coincides with $\mu$ on $\mathcal{R}$
    \item $\mu$ extends to the countably additive extension to the $\sigma$-ring generated by $\mathcal{R}$
\end{enumerate}
If one replaces ring by algebra above, the extension in (4) is unique.
\label{thm:caratheodory}
\end{theorem}
\begin{proof}
See \cite[1.11.8~Theorem]{Bo07:MeasureTheory}.
\end{proof}

Let $X$ be a topological space.
Borel sigma algebra.
Define outer regular, inner regular, locally finite measures.